% Partie : logique du premier ordre
% Chapitre : complétude
% Section : démonstration directe de la compacité

\documentclass[../../../include/open-logic-section]{subfiles}

\begin{document}

\iftag{FOL}
      {\olfileid[fr]{fol}{com}{cpd}}
      {\olfileid[fr]{pl}{com}{cpd}}

\olsection{Une démonstration directe du théorème de compacité}

Nous pouvons démontrer directement le théorème de compacité, sans
invoquer celui de complétude, en reprenant les idées de la
démonstration de ce dernier. Nous étions partis d'un ensemble
cohérent~$\Gamma$ \iftag{FOL}{d'!!{sentence}s}{de !!{formula}s}, l'avions étendu en un ensemble
$\Gamma^*$ cohérent\iftag{FOL}{,
saturé}{} et !!{complete}, puis avions montré que
\iftag{FOL}{le modèle des termes~$\Struct{M(\Gamma^*)}$ satisfait
les !!{sentence}s de~$\Gamma$ sans égalité, tandis que le modèle
quotient permet de traiter les ensembles d'énoncés avec égalité}{la !!{valuation}~
$\pAssign{v(\Gamma^*)}$ satisfait toutes les !!{formula}s
de~$\Gamma$}. Ainsi, $\Gamma$ est satisfaisable.

La même méthode permet de montrer qu'un ensemble \iftag{FOL}{d'!!{sentence}s}{de !!{formula}s}
finiment satisfaisable est satisfaisable. Il suffit de démontrer les
versions correspondantes des résultats menant au lemme de vérité, en
remplaçant « cohérent » par « finiment satisfaisable ».

\begin{prop}
\ollabel{prop:fsat-ccs}
Supposons $\Gamma$ !!{complete} et finiment satisfaisable. Alors :
\begin{enumerate}
\tagitem{prvAnd}{$(!A\land !B)\in\Gamma$ si et seulement si
  $!A\in\Gamma$ et $!B\in\Gamma$.}{}

\tagitem{prvOr}{$(!A\lor !B)\in\Gamma$ si et seulement si
  $!A\in\Gamma$ ou $!B\in\Gamma$.}{}

\tagitem{prvIf}{$(!A\lif !B)\in\Gamma$ si et seulement si
  $!A\notin\Gamma$ ou $!B\in\Gamma$.}{}
\end{enumerate}
\end{prop}

\tagprob{FOL}
\begin{prob}
Démontrez la \olref[fol][com][cpd]{prop:fsat-ccs} sans utiliser~$\Proves$.
\end{prob}
\tagendprob

\tagprob{notFOL}
\begin{prob}
Démontrez la \olref[pl][com][cpd]{prop:fsat-ccs} sans utiliser~$\Proves$.
\end{prob}
\tagendprob

\iftag{FOL}{%
\begin{lem}
\ollabel{lem:fsat-henkin}
Tout ensemble~$\Gamma$ finiment satisfaisable peut être étendu en un
ensemble~$\Gamma'$ saturé et finiment satisfaisable.
\end{lem}
}{}

\tagprob{FOL}
\begin{prob}
Démontrez le \olref[fol][com][cpd]{lem:fsat-henkin}. (Indication :
l'étape décisive consiste à montrer que, si $\Gamma_n$ est finiment
satisfaisable, $\Gamma_n\cup\{!D_n\}$ l'est aussi, sans invoquer les
!!{derivation}s ni la cohérence.)
\end{prob}
\tagendprob

\iftag{FOL}{
\begin{prop}\ollabel{prop:fsat-instances}
Supposons $\Gamma$ complet, finiment satisfaisable et saturé, et
supposons que les variables libres de~$!A(x)$ appartiennent à~$\{x\}$.
\begin{tagenumerate}{prvEx,prvAll}
\tagitem{prvEx}{$\lexists[x][!A(x)]\in\Gamma$ si et seulement si
  $!A(t)\in\Gamma$ pour au moins un terme clos~$t$.}{}
\tagitem{prvAll}{$\lforall[x][!A(x)]\in\Gamma$ si et seulement si
  $!A(t)\in\Gamma$ pour tout terme clos~$t$.}{}
\end{tagenumerate}
\end{prop}
}{}

\tagprob{FOL}
\begin{prob}
Démontrez la \olref[fol][com][cpd]{prop:fsat-instances}.
\end{prob}
\tagendprob

\begin{lem}
\ollabel{lem:fsat-lindenbaum}
Tout ensemble~$\Gamma$ finiment satisfaisable peut être étendu en un
ensemble~$\Gamma^*$ complet et finiment satisfaisable.
\end{lem}

\tagprob{FOL}
\begin{prob}
Démontrez le \olref[fol][com][cpd]{lem:fsat-lindenbaum}. (Indication :
l'étape décisive consiste à montrer que, si $\Gamma_n$ est finiment
satisfaisable, alors $\Gamma_n\cup\{!A_n\}$ ou
$\Gamma_n\cup\{\lnot !A_n\}$ est finiment satisfaisable.)
\end{prob}
\tagendprob

\tagprob{notFOL}
\begin{prob}
Démontrez le \olref[pl][com][cpd]{lem:fsat-lindenbaum}. (Indication :
l'étape décisive consiste à montrer que, si $\Gamma_n$ est finiment
satisfaisable, alors $\Gamma_n\cup\{!A_n\}$ ou
$\Gamma_n\cup\{\lnot !A_n\}$ est finiment satisfaisable.)
\end{prob}
\tagendprob

\begin{thm}[Compacité]
\ollabel{thm:compactness-direct}
$\Gamma$ est satisfaisable si et seulement s'il est finiment
satisfaisable.
\end{thm}

\begin{proof}
Si $\Gamma$ est satisfaisable, il existe
\iftag{FOL}{!!a{structure}~$\Struct{M}$}{!!a{valuation}~$\pAssign{v}$}
telle que $\iftag{FOL}{\Sat{M}}{\pSat{v}}{!A}$ pour \iftag{FOL}{tout !!{sentence}}{toute !!{formula}}
$!A\in\Gamma$. Cette
\iftag{FOL}{$\Struct{M}$}{$\pAssign{v}$} satisfait évidemment aussi
chaque sous-ensemble fini de~$\Gamma$ ; $\Gamma$ est donc finiment
satisfaisable.

Supposons maintenant $\Gamma$ finiment satisfaisable.\iftag{FOL}{
D'après le \olref{lem:fsat-henkin}, il existe un ensemble saturé et
finiment satisfaisable $\Gamma'\supseteq\Gamma$.}{} D'après le
\olref{lem:fsat-lindenbaum}, on peut étendre
$\iftag{FOL}{\Gamma'}{\Gamma}$ en un ensemble~$\Gamma^*$ complet
et finiment satisfaisable\iftag{FOL}{, qui reste saturé}{}.
Les clauses de la \olref[mod]{defn:termmodel} définissent
\iftag{FOL}{le modèle des termes~$\Struct{M(\Gamma^*)}$}
{la !!{valuation}~$\pAssign{v(\Gamma^*)}$}.
\iftag{FOL}{Le langage enrichi contient des constantes. Le modèle
des termes est couvert, et la \olref[mod]{prop:quant-termmodel}
s'applique à ses instances closes.}{}

Si une partie finie de~$\Gamma^*$ entraîne $!A$, alors $!A$ appartient
à~$\Gamma^*$. En effet, si $!A$ en était absent, la complétude de
l'ensemble donnerait $\lnot !A\in\Gamma^*$, en contradiction avec
la satisfaisabilité de cette partie finie augmentée de~$\lnot !A$.
Ainsi, les vérités logiques appartiennent à~$\Gamma^*$ et la
satisfaisabilité de ses parties finies exclut l'absurde. Pour chaque
formule décidée par~$\Gamma^*$, la formule et sa négation ne peuvent
appartenir toutes deux à l'ensemble.
La démonstration du lemme de vérité
(\olref[mod]{lem:truth}) se transpose en remplaçant la
\olref[ccs]{prop:ccs} par la \olref{prop:fsat-ccs}\iftag{FOL}{, et la
\olref[hen]{prop:saturated-instances} par la
\olref{prop:fsat-instances}}{}.
\iftag{FOL}{Pour les énoncés contenant~$\eq$, il faut en outre
passer au quotient du modèle des termes comme dans la section sur
l'égalité. La relation $t\approx t'$ définie par
$\eq[t][t']\in\Gamma^*$ est une congruence : la complétude et la
satisfaisabilité finie excluent chaque violation des lois d'égalité
ou de substitution. Par exemple, deux égalités $\eq[t][t']$ et
$\eq[t'][t'']$ entraînent $\eq[t][t'']$, qui appartient donc à~$\Gamma^*$.
La même propriété de conséquence finie s'applique à la réflexivité,
à la symétrie et au remplacement dans les fonctions et les prédicats.
Dans le quotient, le cas de l'égalité du lemme de vérité revient à
comparer les classes des deux termes clos. Les autres cas utilisent
les mêmes propriétés que dans le modèle des termes. On obtient ainsi,
avec ou sans égalité, une structure dont le réduit au langage initial
satisfait~$\Gamma$.}{Ainsi la valuation satisfait toutes les formules
de~$\Gamma$.}
\end{proof}

\tagprob{FOL}
\begin{prob}
Rédigez la démonstration complète de la version du lemme de vérité
(\olref[fol][com][mod]{lem:truth}) nécessaire à la démonstration du
\olref[fol][com][cpd]{thm:compactness-direct}.
\end{prob}
\tagendprob

\tagprob{notFOL}
\begin{prob}
Rédigez la démonstration complète de la version du lemme de vérité
(\olref[pl][com][mod]{lem:truth}) nécessaire à la démonstration du
\olref[pl][com][cpd]{thm:compactness-direct}.
\end{prob}
\tagendprob

\end{document}
